\section{Discussion}\label{sec:discussion}

Conditional evidence gives a credit analyst a comparison that marginal WoE cannot supply: how the observed characteristic changes default odds after the rest of this record is known. The numerical example's increase from 10 to 20 percent risk corresponds to 0.81 log-odds units. In the simulation, the comparison recovers a deliberately constructed file-age interaction that a single marginal class value misses. It also shows why a small signed average cannot establish that a block is irrelevant. Installment evidence varies across applicants, and removing its direct risk term eliminates its last-entry evidence while leaving most of its order-averaged evidence.

The estimator must be judged against the conditional evidence for each applicant. Averaging estimation errors across applicants assesses how accurately these local quantities are recovered; it does not turn them into a measure of overall model fit. Projection has lower error than separate refits in this simulation, with both smaller reduced-posterior errors and more cancellation of errors shared with the full model. This makes it useful for describing the evidence implied by an accurate model, but does not establish robustness to a misspecified full model. Leaf-local frequencies provide a different, coarser comparison. Concentrating their weights improves on marginal WoE, yet the installment estimate remains less accurate than assigning zero conditional evidence. The amount of recoverable evidence therefore matters as much as the improvement over a weak baseline.

For adverse action, the relevant scored comparison remains the declined applicant against specified passing records. In the featured case, delinquency leads against the broader score band and revolving credit leads against references with more similar delinquency histories. Both comparisons reconcile the score gap. They identify different disadvantages because the reference populations differ. A lender interpreting the leading block must therefore consider the reference characteristics and whether the block description accurately represents the scored factors. Exact allocation does not choose the reference policy or resolve field-level wording within a multi-field block.

The population being described is also consequential. Booked accounts can have a different distribution of an omitted characteristic from through-the-door applicants with the same remaining record. In the selection experiment, booked-population evidence is larger on average for the blocks used by the earlier policy, with revolving and delinquency differences of 0.54 and 0.74 logit units among current-model declines. Those numbers describe the studied policy and population; selection on information outside the observed record would also alter the full-record posterior.

The simulation was designed to expose conditional effects and permits numerical assessment of evidence error. It does not reproduce all bureau missingness, monthly histories or population changes. Its full model is well fitted, so the projection advantage has not been tested under systematic model misspecification or in a generating model without the constructed reversal. The oracle sensitivity checks produce small discrepancies between numerical approximations, but cannot detect bias common to those approximations. Six blocks permit exact order averaging; approximation error with many more blocks is not studied.

HELOC supplies a real-data application under an imposed approval rule, without a population oracle or observed notice reasons. Its calibration, special-value handling and small local samples can affect evidence estimates, but the separate contribution of each was not identified. The attribution hybrids preserve within-block identities but were not checked for cross-block support. Regulatory interpretation is confined to the cited Regulation B text and Official Interpretations; these do not specifically address grouped Shapley attribution. The cited 2022 circular was withdrawn in May 2025 \citep{withdrawn}.

\section{Conclusion}

Conditional weight of evidence measures the adverse or favorable information a characteristic adds for an individual applicant, given the remaining credit record. In this simulation, regression of full-model probabilities estimates that information more accurately than separate outcome-based refits. Order averaging provides a related allocation of total evidence across characteristics, including information they share.

Passing-reference choice changes the leading attributed block for about a third of simulated declines. Interpreting a characteristic's local risk information and explaining its contribution to a score gap therefore involve distinct comparisons. Conditional WoE supplies the first; the second requires a specified passing population and an allocation of the fitted score difference.
